\section{Method}
\label{sec:method}

This section develops notation, the generative model, the inference scheme, and the theoretical results in turn.

\subsection{Notation}
\label{sec:notation}

We collect the principal symbols used throughout the paper in Table~\ref{tab:notation}. The model is symbol-heavy and a notation index materially aids review.

\begin{table}[H]
\centering
\caption{Principal notation used throughout the paper.}
\label{tab:notation}
\small
\setlength{\tabcolsep}{4pt}
\begin{tabularx}{\textwidth}{lY}
\toprule
Symbol & Meaning \\
\midrule
$n$ & number of features \\
$K$ & number of replicates \\
$X_{i,j}$ & raw score of feature $i$ on replicate $j$ \\
$\bU_i=(U_{i,1},\ldots,U_{i,K})$ & uniformized vector with $U_{i,j}=F_j(X_{i,j})$ \\
$Z_i\in\{0,1\}$ & latent reproducibility class indicator \\
$\pi=\Pr(Z_i=1)$ & reproducible-component proportion \\
$c_0,c_1$ & irreproducible and reproducible-component copula densities \\
$F_j$ & marginal CDF of replicate $j$ \\
$\sigma\in[0,1)$ & Pitman--Yor discount parameter \\
$\alpha>0$ & Pitman--Yor concentration parameter \\
$T,T_n$ & variational truncation level or theoretical sieve truncation level \\
$T_{-i}$ & number of occupied reproducible clusters after removing feature $i$ \\
$n_{m,-i}$ & count in occupied cluster $m$ after removing feature $i$ \\
$\theta_m,t_m$ & continuous parameters and atomic-type indicator of atom $m$ \\
$w_m,V_m$ & PY mixture weight and stick-breaking ratio \\
$z_i\in\mathbb N$ & cluster assignment within the reproducible component \\
$M_j$ & Bernstein basis order for marginal $j$ \\
$S_{i,j}$ & Bernstein basis-membership indicator for score $X_{i,j}$ \\
$H$ & number of auxiliary atoms in Algorithm~\ref{alg:mcmc} \\
$h(\cdot,\cdot)$ & Hellinger distance \\
$\epsilon_n,\delta_n$ & posterior contraction rate and local-idr error scale \\
$\widehat k_\alpha$ & number of rejections selected by the Bayes-mFDR step-up rule \\
\bottomrule
\end{tabularx}
\end{table}


\subsection{Generative model}
\label{sec:model}

\subsubsection{Latent class and copula representation}

For $i = 1, \ldots, n$, let $Z_i \in \{0, 1\}$ be the latent reproducibility class with $\Pr(Z_i = 1) = \pi$, where $Z_i = 1$ denotes the reproducible component and $Z_i = 0$ the irreproducible component. The population uniformized vector is $\bU_i = (F_1(X_{i,1}), \ldots, F_K(X_{i,K}))^\top \in [0,1]^K$, where $F_j$ is the marginal CDF of replicate $j$; in practice $\bU_i$ is replaced by its empirical scaled-rank counterpart prior to the marginal-model step. Conditional on $Z_i = z$, the uniformized vector has copula density
\begin{equation}
\bU_i \mid Z_i = z\ \sim\ c_z(\bU_i; \btheta_z),
\label{eq:cop-density}
\end{equation}
where, by the IDR convention, $c_0$ is the independence copula on $[0,1]^K$ and $c_1$ is the reproducible-component copula. We retain this convention throughout for compatibility with the IDR literature; alternative weakly-dependent choices for $c_0$ are possible and would not affect the structural results. The local idr is the per-feature posterior of the irreproducible class, $\idr_i = \Pr(Z_i = 0 \mid \bU_i)$, and the aggregated $\IDR(\gamma)$ is the marginal FDR (mFDR) of the rejection set $\mathcal{R}(\gamma) = \{i : \idr_i < \gamma\}$,
\begin{equation}
\IDR(\gamma) \;=\; \frac{\E\!\left[\sum_{i=1}^n (1-Z_i)\,\bm{1}\{\idr_i < \gamma\}\right]}{\E\!\left[\sum_{i=1}^n \bm{1}\{\idr_i < \gamma\}\right]},
\label{eq:agg-idr}
\end{equation}
which serves as the threshold-free analogue of the BH cutoff and is the operational target of the Sun--Cai-type step-up rule \citep{SunCai2007}.

\subsubsection{Reproducible-component copula}

Conditional on $Z_i = 1$, the copula density is the Pitman--Yor mixture
\begin{equation}
c_1(\bU_i) \;=\; \sum_{m=1}^\infty w_m\, c_{\theta_m}(\bU_i),\qquad (w_1, w_2, \ldots) \sim \GEM(\alpha, \sigma), \quad (\theta_m, t_m) \stackrel{\mathrm{iid}}{\sim} P_0,
\label{eq:py-mix}
\end{equation}
where each atom is a pair $(\theta_m, t_m)$ comprising continuous parameters $\theta_m$ and a discrete atomic-type indicator $t_m \in \{\text{Gauss}, t_5, \text{Clayton}, \text{Gumbel}\}$ that selects the copula family. The base measure $P_0$ factorizes as $P_0(\theta, t) = P_0(\theta \mid t)\, P_0(t)$, with $P_0(t)$ uniform on the four families. For Gaussian atoms, the structural class of the correlation matrix $R(\theta_m)$ is itself drawn at cluster instantiation from a uniform prior on the three classes \emph{exchangeable} ($R = (1-\rho) I + \rho \one\one^\top$), \emph{one-factor} ($R = D_1 + \lambda_1 \lambda_1^\top$), and \emph{two-factor} ($R = D_2 + \Lambda \Lambda^\top$ with $\Lambda \in \R^{K \times 2}$), where $D_q$ are diagonal matrices chosen so $R$ has unit diagonal. We exclude vine constructions for analytical tractability and defer them to future work. For non-Gaussian atomic types, $\theta_m$ comprises the corresponding family-specific dependence parameters.

\subsubsection{Replicate-specific marginals}

Per-replicate marginals $F_j$ have Bernstein--Dirichlet priors,
\begin{equation}
F_j(x) \;=\; \sum_{r=1}^{M_j} p_{j,r}\, B_r\!\left(F_j^{(0)}(x);\, M_j\right),\qquad B_r(t; M) \;=\; I_t(r, M-r+1),\qquad \bm{p}_j \sim \Dir(\alpha_F \one_{M_j}),
\label{eq:bernstein}
\end{equation}
where $B_r(t; M)$ is the cumulative Bernstein basis of order $M_j$ at index $r$, identified with the regularized incomplete beta function $I_t(r, M-r+1)$; $F_j^{(0)}$ is a coarse pilot CDF (e.g., a kernel-smoothed empirical CDF with a wide bandwidth); and $\alpha_F > 0$ is a smoothing hyperparameter. The Bernstein--Dirichlet construction is canonical for densities on $[0,1]$ \citep{Petrone1999,PetroneWasserman2002}.

\subsubsection{Identifiability constraint}

The reproducible-component mixture is constrained to atoms with positive lower-orthant dependence (PLOD): for every atom index $m$,
\begin{equation}
C_{\theta_m}(u, \ldots, u) \;>\; u^K \quad \text{for all } u \in (0, 1),
\label{eq:plod}
\end{equation}
where $C_{\theta_m}$ is the copula function corresponding to density $c_{\theta_m}$. The PLOD condition \citep[Definition 2.1]{Joe1997} is satisfied by Gaussian copulas with all pairwise correlations strictly positive, by $t$ copulas with positive correlation, and by Clayton and Gumbel copulas. Unlike the strict lower-tail-dependence condition $\lambda_L > 0$, PLOD does not exclude Gaussian atoms with $|\rho_{jk}| < 1$ and is therefore compatible with the dominant Gaussian-copula base measure. The role of \eqref{eq:plod} is to break the label symmetry between the reproducible and irreproducible components: the independence copula $c_0$ satisfies $C_0(u,\dots,u) = u^K$ exactly, so any atom with strict PLOD is distinguishable from $c_0$.

\subsubsection{Hyperpriors}
\label{sec:hyperpriors}

We place $\pi \sim \Beta(1,1)$, $\alpha \sim \mathrm{Gamma}(1,1)$, and $\sigma \sim \mathrm{Uniform}[0,1)$, the last specification being equivalent to $\Beta(1,1)$ truncated to exclude the boundary point $\sigma = 1$ excluded by the PY definition. The Bernstein smoothing hyperparameter satisfies $\alpha_F \sim \mathrm{Gamma}(1,1)$. The Bernstein degree $M_j$ has a discrete prior on $\{5, 10, 20, 40, 80\}$ with weights proportional to $1/M_j$, normalized so that
\begin{equation}
\Pr(M_j = m) \;=\; \frac{1/m}{\sum_{m' \in \{5,10,20,40,80\}} 1/m'}, \qquad m \in \{5, 10, 20, 40, 80\}.
\label{eq:m-prior}
\end{equation}
The base measure $P_0(\theta \mid t)$ on Gaussian-copula structures uses an LKJ($\eta = 2$) prior \citep{LewandowskiKurowickaJoe2009} on the correlation matrix $R$, with the induced prior on the Cholesky factor $L$ following from $R = LL^\top$. The choice $\eta = 2$ provides mild concentration toward the identity matrix and serves as a weakly informative default that biases toward simpler dependence structures in the absence of strong evidence. The structural-class indicator and the atomic-type indicator each take uniform priors over their respective four-element and three-element supports.

\subsection{Inference}
\label{sec:inference}

\subsubsection{Slice-sampled P\'olya-urn MCMC}

We employ a slice-sampled P\'olya-urn MCMC scheme based on Neal's Algorithm~8 \citep{Neal2000}, extended from the Dirichlet process to the Pitman--Yor process via the predictive of \citet{Pitman1996BlackwellMacQueen} and \citet{IshwaranJames2001}. Algorithm~\ref{alg:mcmc} details one full sweep of the sampler. The conditional posterior class probability used in Step~1 is
\begin{equation}
\widehat\pi_i \;=\; \Pr(Z_i = 1 \mid \bU_i, \btheta, \bw, \bm{F}) \;=\; \frac{\pi\, c_1(\bU_i; \btheta, \bw)}{\pi\, c_1(\bU_i; \btheta, \bw) + (1 - \pi)\, c_0(\bU_i)},
\label{eq:cond-class}
\end{equation}
with $c_1(\bU_i; \btheta, \bw) = \sum_m w_m\, c_{\theta_m}(\bU_i)$ and $c_0(\bU_i) = 1$ for the independence copula. The atomic-type indicator update in Step~3 uses an independence Metropolis proposal that draws the candidate type uniformly from the four families. The Bernstein-weight update in Step~4 is conjugate after introducing latent basis-membership indicators $S_{i,j} \in \{1, \ldots, M_j\}$ with $\Pr(S_{i,j} = r \mid \cdot) \propto p_{j,r}\, b_r(F_j^{(0)}(X_{i,j}); M_j)$, where $b_r(t; M) = M\binom{M-1}{r-1}t^{r-1}(1-t)^{M-r}$ is the Bernstein density basis; the conditional posterior is then $\bm{p}_j \mid \cdot \sim \Dir(\alpha_F + n_{j,1}, \ldots, \alpha_F + n_{j,M_j})$ with $n_{j,r} = \#\{i : S_{i,j} = r\}$. The PY hyperparameter updates use the auxiliary-variable scheme of \citet{EscobarWest1995} for $\alpha$ and a slice-sampled update for $\sigma$ on $[0,1)$.

\begin{algorithm}[H]
\DontPrintSemicolon
\caption{One sweep of slice-sampled P\'olya-urn MCMC for PY-IDR.}
\label{alg:mcmc}
\KwIn{Current state $(\pi, \btheta, \bw, \bm{F}, \alpha, \sigma)$, $H$ auxiliary atoms, observed scores $\bX_{1:n}$.}
\KwOut{Updated state.}
\textbf{Class assignment.} For each $i$, draw $Z_i \sim \mathrm{Bernoulli}(\widehat\pi_i)$ via Eq.~\eqref{eq:cond-class}\;
\textbf{Cluster reassignment.} For each $i$ with $Z_i = 1$, sample cluster index $z_i$ via Algorithm~8 with $H$ auxiliary atoms drawn from $P_0$\;
\textbf{Atom-parameter update.} For each occupied cluster $m$, update $\theta_m$ via NUTS on the unconstrained Cholesky parameterization $R(\theta_m) = LL^\top$ with $L_{j,j} = \exp(\eta_{j,j})$ and $L_{j,k}$ free for $j > k$\;
\textbf{Atomic-type update.} For each occupied cluster $m$, update $t_m$ by independence Metropolis with uniform proposal over the four atomic families\;
\textbf{Marginal-weight update.} Sample latent basis-membership indicators $S_{i,j}$ and update $\bm{p}_j \sim \Dir(\alpha_F + n_{j,1}, \ldots, \alpha_F + n_{j,M_j})$\;
\textbf{Marginal-degree update.} Update $(M_j, \alpha_F)$ by Metropolis-within-Gibbs with random-walk proposals on $\log \alpha_F$ and a discrete proposal on $M_j$ over $\{5, 10, 20, 40, 80\}$\;
\textbf{PY hyperparameter update.} Update $\alpha$ by the auxiliary-variable scheme of \citet{EscobarWest1995}; update $\sigma \in [0,1)$ by slice sampling\;
\textbf{Mixing-proportion update.} Sample $\pi \sim \Beta(1 + n_1, 1 + n_0)$ where $n_z = \#\{i : Z_i = z\}$\;
\Return{Updated state.}\;
\end{algorithm}

\paragraph{Initialization and diagnostics.}
We initialize the sampler from EM-IDR estimates of $(\pi, c_1)$ on the bivariate marginals collapsed to $K = 2$ pairs, with cluster assignments seeded by a $k$-means split of the rank-transformed data into a small number of groups. We monitor convergence by the split-$\widehat R$ statistic with threshold $1.01$ and require an effective sample size of at least $400$ per parameter before reporting posterior summaries.

\paragraph{Complexity.}
Each MCMC sweep is $O(nKT + K^3 T)$, where $T$ is the number of active clusters and $K$ is the replicate count. The first term is dominated by the per-cluster, per-feature copula-density evaluation; the second arises from the Cholesky factorization step on each cluster's $K \times K$ correlation matrix.

\subsubsection{Variational alternative}
\label{sec:vi}

A structured mean-field family with truncation $T$ at the PY stick yields a variational ELBO that we optimize by stochastic natural-gradient ascent in the framework of \citet{HoffmanBleiWangPaisley2013}, with consistency guarantees for the truncated parametric model following \citet{WangBlei2019} and for the nonparametric limit following \citet{AlquierRidgway2020} and \citet{YangPatiBhattacharya2020}. The variational family factorizes as
\[
q(Z, z, \theta, t, w, \pi, F, \alpha, \sigma) \;=\; \prod_i q(Z_i)\,\prod_i q(z_i)\,\prod_{m \leq T} q(\theta_m)\,\prod_{m \leq T} q(t_m)\,q(w)\,q(\pi)\,\prod_j q(F_j)\,q(\alpha)\,q(\sigma),
\]
with $q(Z_i)$ and $q(z_i)$ categorical, $q(\theta_m)$ Gaussian on the unconstrained Cholesky entries, $q(t_m)$ categorical over the four atomic types, $q(w)$ a product of independent $\Beta$ factors on the stick-breaking ratios, $q(\pi)$ a $\Beta$ distribution, $q(F_j)$ a Dirichlet on the Bernstein weights, and $q(\alpha)$, $q(\sigma)$ Gamma and truncated-Beta respectively. Per-iteration cost is $O(nKT \bar d)$ with $\bar d$ the mean atomic-parameter dimension; mini-batches of size $B \in \{500, 2000\}$ are used in practice. As a practical recommendation, we suggest the MCMC scheme for $n \lesssim 10^4$ where exact posterior summaries with credible bands are required, and the variational alternative for $n \gtrsim 10^5$ where the per-sweep cost of MCMC becomes prohibitive and approximate posterior summaries suffice.

\subsubsection{Local idr, aggregated IDR, and Bayes-mFDR}

The posterior local idr is the marginal posterior probability of the irreproducible class,
\begin{equation}
\widehat{\idr}_i \;=\; \Pr(Z_i = 0 \mid \mathrm{data}) \;=\; \int \Pr(Z_i = 0 \mid \mathrm{data}, \btheta, \bw, \bm{F})\, \pi(\btheta, \bw, \bm{F} \mid \mathrm{data})\, d(\btheta, \bw, \bm{F}).
\label{eq:idr}
\end{equation}
The Bayes marginal false discovery rate at threshold $\gamma$ is
\begin{equation}
\mathrm{Bayes\text{-}mFDR}(\gamma) \;=\; \frac{\E[V(\gamma) \mid \mathrm{data}]}{\E[R(\gamma) \mid \mathrm{data}]},\qquad V(\gamma) = \sum_i (1 - Z_i)\bm{1}\{\widehat{\idr}_i < \gamma\},\;\; R(\gamma) = \sum_i \bm{1}\{\widehat{\idr}_i < \gamma\},
\label{eq:bayes-mfdr}
\end{equation}
where $V$ and $R$ are the false-rejection and total-rejection counts under threshold $\gamma$. The Sun--Cai-type step-up cutoff is the largest $\gamma_\alpha$ such that the running average of sorted $\widehat{\idr}_{(1)} \leq \cdots \leq \widehat{\idr}_{(n)}$ satisfies $k^{-1} \sum_{j=1}^k \widehat{\idr}_{(j)} \leq \alpha$ at $k = R(\gamma_\alpha)$ \citep{SunCai2007}.

\begin{remark}
The local idr is a per-feature posterior used for ranking, whereas the aggregated IDR is a rejection-set average used for thresholding. The two are connected through the running-average threshold rule above, which is the empirical-Bayes plug-in counterpart of the optimal compound-decision rule of \citet{SunCai2007}.
\end{remark}

\subsection{Theoretical properties}
\label{sec:theory}

We state four results on identifiability, posterior contraction, asymptotic mFDR control, and ergodicity of the sampler. Full proofs and supporting lemmas appear in the appendix.

\subsubsection{Assumptions}

\begin{assumption}[Smoothness and boundedness]
\label{ass:smooth}
The true reproducible-component copula density $c_1^*$ is $\beta$-H\"older with $\beta \geq 1$ on $[0,1]^K$ and is bounded above by a constant $B < \infty$, $\sup_{u \in [0,1]^K} c_1^*(u) \leq B$. The true marginal densities $f_j^*$ are $\beta$-H\"older on the support of $X_{\cdot, j}$.
\end{assumption}

\begin{assumption}[Reproducible mass]
\label{ass:pi}
$\pi^* \in (0, 1)$ is bounded away from $\{0, 1\}$.
\end{assumption}

\begin{assumption}[Prior KL-support]
\label{ass:kl}
For every $\epsilon > 0$, the prior on the reproducible-component copula satisfies $\Pi(\KL(c_1^*, c_1) < \epsilon^2) > 0$, and analogously for each marginal.
\end{assumption}

\begin{assumption}[PY truncation]
\label{ass:py}
The truncation level satisfies $T_n \asymp n \epsilon_n^2 / \log n$ as $n \to \infty$, where $\epsilon_n$ is the contraction rate of Theorem~\ref{thm:contract}.
\end{assumption}

\begin{assumption}[Slice-sampler regularity]
\label{ass:slice}
The transition kernel of the slice-sampled P\'olya-urn sampler admits a positive transition density on a small set $C \subset \Theta$ with positive Lebesgue measure $\mu(C) > 0$, and a Lyapunov function $V(\btheta) = 1 + \norm{\btheta}_F^2$ (with $\norm{\cdot}_F$ the Frobenius norm on the matrix-valued components of $\btheta$) satisfies the geometric drift condition $\E[V(\btheta') \mid \btheta] \leq \lambda V(\btheta) + b\,\bm{1}_C(\btheta)$ for some $\lambda < 1$, $b < \infty$.
\end{assumption}

\subsubsection{Identifiability}

\begin{theorem}[Identifiability]
\label{thm:id}
Under Assumptions~\ref{ass:smooth} and~\ref{ass:pi}, the joint $(\pi, c_0, c_1, F_1, \ldots, F_K)$ is identifiable up to label permutation of mixture atoms in $c_1$, given the PLOD constraint~\eqref{eq:plod} on the reproducible-component atoms.
\end{theorem}

\begin{proof}
The argument proceeds in three steps.

\emph{Step 1 (separation of marginals).} By Sklar's theorem \citep{Sklar1959}, the marginals $\{F_j\}_{j=1}^K$ are uniquely determined from the joint distribution of $\bX$, since each $F_j$ is the marginal CDF of $X_{\cdot,j}$. Substituting $\bU = (F_1(X_1), \ldots, F_K(X_K))$ reduces the identification problem to that of the copula mixture $\pi c_1 + (1-\pi) c_0$ on $[0,1]^K$.

\emph{Step 2 (separation of components via PLOD).} The independence copula $c_0$ satisfies $C_0(u, \ldots, u) = u^K$ for all $u \in (0,1)$, while every PLOD atom in the reproducible component satisfies $C(u, \ldots, u) > u^K$. The diagonal section
\[
\Delta(u) \;\myeq\; C_{\mathrm{joint}}(u, \ldots, u) - u^K \;=\; \pi \int (C_\theta(u, \ldots, u) - u^K)\, dW(\theta)
\]
where $W$ is the implied mixing measure on atoms, is identically zero only when $\pi = 0$ and is strictly positive when $\pi > 0$. Differentiating $\Delta$ in $u$ at the boundary recovers $\pi$ uniquely as the residual mass; the irreproducible component is then identified as the unique copula with diagonal section $u^K$.

\emph{Step 3 (identification of the reproducible mixture).} Identification of $c_1$ from the reproducible mixture follows from the linear-independence criterion of \citet{YakowitzSpragins1968}: a finite mixture of densities is identifiable up to label permutation if and only if the candidate component densities form a linearly independent family. Gaussian copula densities with distinct correlation matrices are real-analytic in $\bU$ on the open cube $(0,1)^K$, and any nontrivial linear combination of them with distinct correlation matrices cannot vanish on a set of positive Lebesgue measure, since real-analytic functions vanishing on a positive-measure set must be identically zero. The heavy-tail families ($t_5$, Clayton, Gumbel) are likewise linearly independent of each other and of the Gaussian family. Hence the Yakowitz--Spragins hypothesis is satisfied, and $c_1$ is identified up to the residual permutation of atom labels.
\end{proof}

\begin{remark}
\label{rem:amr-vs-yakowitz}
We use \citet{YakowitzSpragins1968} rather than \citet{AllmanMatiasRhodes2009} as the primary identifiability tool. The latter's Kruskal-based theorem requires conditional independence of the observed variables given the latent class, which is violated by Gaussian-copula mixture components since the variates within each component are correlated. The Yakowitz--Spragins linear-independence criterion is well-suited to copula-mixture identifiability and avoids this structural obstacle.
\end{remark}

\subsubsection{Posterior contraction}

\begin{theorem}[Hellinger posterior contraction]
\label{thm:contract}
Under Assumptions~\ref{ass:smooth}--\ref{ass:py}, the posterior contracts in Hellinger distance at rate $\epsilon_n = n^{-\beta/(2\beta+K)} (\log n)^t$ for some $t > 0$:
\[
\Pi\!\left(h(f, f^*) > M \epsilon_n \mid \bX_{1:n}\right) \;\to\; 0 \qquad \text{in } P^*\text{-probability,}
\]
for sufficiently large $M$, where $f$ is the joint density of $\bX_i$ implied by the model and $f^*$ is the truth.
\end{theorem}

\begin{proof}[Proof outline]
We invoke the master contraction theorem of \citet{GhosalvdV2017}, Theorem~8.9. Three conditions are verified.

\emph{Prior-mass condition.} For any $c_1^* \in C^\beta([0,1]^K)$ satisfying Assumption~\ref{ass:smooth}, an extension of the multivariate Dirichlet-mixture contraction theorem of \citet{ShenTokdarGhosal2013} from densities on $\R^K$ to copulas on $[0,1]^K$, together with the Pitman--Yor adaptation of \citet{Scricciolo2014BAPY}, shows that there exists a mixture with $T_n = O(n \epsilon_n^2 / \log n)$ Gaussian-copula atoms within $\epsilon_n$ in KL divergence of $c_1^*$. The transformation to the cube introduces an additional Sklar-based regularity step that we discharge in Lemma~\ref{lem:sklar-regularity} of the appendix. The Bernstein--Dirichlet prior on each marginal places positive mass on KL balls of radius $\epsilon_n$ around any $\beta$-H\"older marginal density, by the adaptive contraction analysis of \citet{Rousseau2010BetaMix}; the slower non-adaptive rate of \citet{Ghosal2001Bernstein} and the foundational consistency results of \citet{Petrone1999} and \citet{PetroneWasserman2002} are background. The KL prior-mass condition then follows from the mixture-entropy decomposition of Lemma~\ref{lem:mixent} in the appendix, which gives $\KL(f, f^*) \lesssim \KL(c, c^*) + \sum_j \KL(f_j, f_j^*)$ when $\sup c, \sup c^* \leq B < \infty$.

\emph{Sieve construction.} The sieve $\mathcal{F}_n = \{\sum_{m \leq T_n} w_m\, c_{(\theta_m, t_m)}(\cdot) : \norm{\theta_m}_F \leq M_n\}$ has covering numbers bounded by Lemma~\ref{lem:pytruncation} of the appendix: $\log N(\epsilon, \mathcal{F}_n, h) \leq C T_n K^2 \log(M_n/\epsilon) + C T_n$, which is $\exp(n \epsilon_n^2)$ for $M_n = n^a$, $a < \infty$.

\emph{Sieve complement.} The PY tail bound under Assumption~\ref{ass:py} gives $\Pi(\mathcal{F}_n^c) = o(\exp(-2n \epsilon_n^2))$, completing the verification.

Combining marginal and copula contraction rates via the mixture-entropy decomposition yields the stated rate $\epsilon_n = n^{-\beta/(2\beta+K)} (\log n)^t$, where the $K$ in the denominator reflects the dimension of the copula and the corresponding curse-of-dimensionality penalty.
\end{proof}

\begin{corollary}[Gaussian-only base measure]
\label{cor:gauss-only}
If the base measure $P_0$ is restricted to Gaussian-copula atoms only, the conclusion of Theorem~\ref{thm:contract} holds with the same rate $\epsilon_n$.
\end{corollary}

\subsubsection{Asymptotic Bayes-mFDR control}

\begin{theorem}[Asymptotic Bayes-mFDR control]
\label{thm:fdr}
Suppose Theorem~\ref{thm:contract} holds, and additionally suppose that the limiting CDF of $\widehat{\idr}_i$ under $P^*$ is continuous and Lipschitz in a neighborhood of the Sun--Cai threshold $\gamma_\alpha^* = \inf\{\gamma : \mFDR^*(\gamma) \leq \alpha\}$. Then the Sun--Cai-type step-up rule applied to the posterior $\widehat{\idr}_i$ controls Bayes-mFDR at the nominal level $\alpha + O(\epsilon_n)$ as $n \to \infty$.
\end{theorem}

\begin{proof}[Proof outline]
Theorem~\ref{thm:contract} implies uniform $L^1$-consistency of the posterior local idr for the oracle local idr $\idr_i^* = \Pr_{P^*}(Z_i = 0 \mid \bU_i)$, since $|\widehat{\idr}_i - \idr_i^*| \leq C h(f, f^*)$ on the event of contraction. The continuity and Lipschitz assumptions on the limiting CDF combined with the continuous mapping theorem applied to the threshold map $\gamma \mapsto \mFDR(\gamma)$ transfer exact mFDR control of the oracle Sun--Cai rule \citep{SunCai2007} to its plug-in counterpart, with error term $O(\epsilon_n)$. The rate-to-mFDR bridge underlying this transfer is the empirical-Bayes multiple-testing analysis of \citet{HellerRosset2021} and the spike-and-slab framework of \citet{CastilloRoquain2020}.
\end{proof}

\subsubsection{Ergodicity of the sampler}

\begin{theorem}[Harris ergodicity, with geometric-rate conjecture]
\label{thm:erg}
Under Assumption~\ref{ass:slice}, the slice-sampled P\'olya-urn kernel $P$ is Harris ergodic. Under the additional hypothesis that the small set $C$ contains the support of the cluster-occupancy configuration with high posterior probability, $P$ is $V$-uniformly geometrically ergodic with rate $\rho < 1$:
\[
\norm{P^k(\btheta, \cdot) - \Pi(\cdot \mid \mathrm{data})}_V \;\leq\; C\, V(\btheta)\, \rho^k \qquad \text{for all } k \geq 0.
\]
\end{theorem}

\begin{proof}[Proof outline]
Harris ergodicity follows from Theorem~13.0.1 of \citet{MeynTweedie2009}: the sampler is irreducible and aperiodic by construction (the auxiliary-atom step of Algorithm~8 supplies positive transition probability between any two configurations differing by one cluster reassignment), and the existence of a small set follows from the positive transition density required by Assumption~\ref{ass:slice}. The geometric-rate strengthening invokes Theorem~16.0.1 of \citet{MeynTweedie2009}, which requires the geometric drift condition of Assumption~\ref{ass:slice} together with the small-set condition; the latter is well-established in finite-dimensional Cholesky parameterizations \citep{RobertsRosenthal2004}, but its extension to the infinite-dimensional configuration space of the PY mixture is a separate technical step that we conjecture but do not establish in full. The recent quantitative-mixing analysis of \citet{Ascolani2024Mixing} for related Dirichlet-process Gibbs samplers provides analogous bounds that support the conjecture.
\end{proof}

\begin{remark}
\label{rem:erg-gap}
We deliberately separate Harris ergodicity from geometric ergodicity in Theorem~\ref{thm:erg}. Harris ergodicity is established and provides asymptotic correctness of MCMC averages; geometric ergodicity is a stronger quantitative-mixing statement that we treat as a conjecture pending a full Foster--Lyapunov drift analysis on the infinite-dimensional configuration space of the PY mixture. We do not appeal to Meyn--Tweedie 16.0.1 as a black box.
\end{remark}

\begin{theorem}[Variational consistency]
\label{thm:vi}
Under Assumptions~\ref{ass:smooth}--\ref{ass:py} and a finite truncation of the PY stick at level $T$, the variational posterior $q^*$ minimizing the KL divergence to the true posterior within the structured mean-field family of Section~\ref{sec:vi} satisfies $h(\widehat f^{q^*}, f^*) \to 0$ in $P^*$-probability. Under additional regularity matching the conditions of \citet{AlquierRidgway2020} for nonparametric BNP models, the rate of convergence is the same $\epsilon_n$ as in Theorem~\ref{thm:contract}.
\end{theorem}

\begin{proof}[Proof outline]
For finite truncation, the parametric Bernstein--von Mises framework of \citet{WangBlei2019} applies and gives the result. For nonparametric extension, the argument follows \citet{AlquierRidgway2020} and \citet{YangPatiBhattacharya2020}: variational consistency holds whenever the true posterior contracts at rate $\epsilon_n$ and the variational family contains a sequence of distributions whose KL divergence to the true posterior is $o(n \epsilon_n^2)$. The structured mean-field family of Section~\ref{sec:vi} has full-rank factors over each parameter block, and standard arguments verify that the projected truth lies within $o(n \epsilon_n^2)$ KL distance of the variational family.
\end{proof}
